\documentclass[a4paper, twoside]{report}

\usepackage[british]{babel}
\usepackage[useregional]{datetime2}
\DTMlangsetup[en-GB]{showdayofmonth=false} %to remove day from date
\usepackage[utf8x]{inputenc}
\usepackage[T1]{fontenc}
\usepackage{listings}
\usepackage{hyperref}
\hypersetup{colorlinks=false}
\usepackage{lscape} %to put things landscape
\usepackage{rotating} %to put things landscape, again
\usepackage{pdflscape} %trying landscape again
\usepackage{subfigure}
\usepackage{amsmath} %for some maths stuff! 
\usepackage{graphicx}
\usepackage[colorinlistoftodos]{todonotes}
\usepackage{csvsimple}
\usepackage{lscape}%for landscape table in sup mat
\usepackage{apalike} %for round brackets in references
\usepackage{titlesec} %for hiding chapter title
%\usepackage{natbib}
\usepackage{gensymb} %degree symbol
\titleformat{\chapter}[display]
  {\normalfont\bfseries}{}{0pt}{\Huge}

%%%%%This stuff is for generating tables%%%%
\usepackage{booktabs} % For \toprule, \midrule and \bottomrule
\usepackage{siunitx} % Formats the units and values
\usepackage{pgfplotstable} % Generates table from .csv

% Setup siunitx:
\sisetup{
  round-mode          = places, % Rounds numbers
  round-precision     = 2, % to 2 places
}
%%%%%%%%%%%%%%%%%%%%%%%%%%%

%% Sets page size and margins
\usepackage[a4paper,top=3cm,bottom=2cm,left=3cm,right=3cm,marginparwidth=1.75cm]{geometry}

\title{Integrating microbial dormancy with biogeography in the changing Arctic}
\author{Amy Solman}
% Update supervisor and other title stuff in title/title.tex

\begin{document}
\begin{titlepage}

\newcommand{\HRule}{\rule{\linewidth}{0.5mm}} % Defines a new command for the horizontal lines, change thickness here

%----------------------------------------------------------------------------------------
%	LOGO SECTION
%----------------------------------------------------------------------------------------

%\includegraphics[width=8cm]{logo.eps}\\[1cm] % Include a department/university logo - this will require the graphicx package
 
%----------------------------------------------------------------------------------------

\center % Center everything on the page

%----------------------------------------------------------------------------------------
%	HEADING SECTIONS
%----------------------------------------------------------------------------------------
\quad\\[1.5cm]
%\textsc{\LARGE MSc Thesis}\\[1.5cm] % Name of your university/college
\textsc{\Large Queen Mary University London \& the Natural History Museum}\\[0.5cm] % Major heading such as course name
%\textsc{\large}\\[0.5cm] % Minor heading such as course title

%----------------------------------------------------------------------------------------
%	TITLE SECTION
%----------------------------------------------------------------------------------------
\makeatletter
\HRule \\[0.4cm]
{ \huge \bfseries \@title}\\[0.4cm] % Title of your document
{ \huge \bfseries \@author}\\[0.4cm]
\HRule \\[1.5cm]
 
%----------------------------------------------------------------------------------------
%	AUTHOR SECTION
%----------------------------------------------------------------------------------------

\begin{minipage}{0.4\textwidth}
\begin{flushleft} \large
\emph{NHM supervisor:}\\
Dr. Anne Jungblut\\
\end{flushleft}
\end{minipage}
~
\begin{minipage}{0.4\textwidth}
\begin{flushright} \large
\emph{QMUL supervisors:} \\
Dr. James Bradley\\
Prof. Kate Heppell\\
% Uncomment the following lines if there's a co-supervisor
%\\[1.2em] % Supervisor's Name
%\emph{Co-Supervisor:} \\
%Dr. Adam Smith % second marker's name
\end{flushright}
\end{minipage}\\[3cm]
\makeatother


%----------------------------------------------------------------------------------------
%	DATE SECTION
%----------------------------------------------------------------------------------------
{\large \today}\\[2cm] % Date, change the \today to a set date if you want to be precise
{\large A research proposal for the QMUL Principal's Studentship funded }\\[0.5cm]
{\large {PhD in Physical Geography}}\\[0.5cm]


\vfill % Fill the rest of the page with whitespace

\end{titlepage}


\renewcommand{\abstractname}{Abstract}
\begin{abstract}
1.  Microorganisms play a fundamental role in nutrient cycling, biogeochemical processes and ecosystem regulation in polar regions. The Arctic is experiencing increased warming compared to the global average, with profound effects for microbial species and their distributions. This study seeks to integrate two understudied areas: microbial biogeography and dormancy in Arctic glaciers and soils. \\

\noindent 2. Dormancy is a survival strategy employed by microorganisms in extreme environments. Active cells can make up less than 10\% of total microbial biomass. Differentiating between dead, dormant and active cells in phylogenetic analysis is integral for understanding community composition and the biogeochemical processes they drive. \\

\noindent 3. Despite their importance, relatively little is known about the biogeographic patterns of microbial communities. Traditionally, microorganisms were assumed to be so abundant, small and free-moving as to be omnipresent and limited only by environmental factors such as pH. However, recent evidence suggests they are subject to many of the same stochastic and deterministic biogeographical mechanisms as macroorganism.\\

\noindent 4. In this study I will incorporate fieldwork, novel laboratory methods and a predictive numerical model, to explore the relationship between microbial dormancy and biogeography in the Arctic across changing seasons and into the future. \\

\noindent \textbf{Keywords:} taxa-area relationships, microbial biogeography, dormancy, niche structure, colonisation-extinction balance, biogeography, climate change, Arctic 
\end{abstract}

\tableofcontents

\chapter{Introduction}

Microorganisms are ubiquitous in the biosphere, contributing significantly to nutrient and carbon cycling. Biological processes in glaciers and ice sheets are dominated almost exclusively by microbial communities \cite{AnesioAm2017Tmog}. Arctic microbes contribute to the build-up of biomass and macronutrients such as carbon, phosphorus and nitrogen in soils \cite{BradleyJamesA2014Mcdi}, as well as impacting the fertilization and productivity of Polar fjords and oceans. They can also modify the albedo of ice surfaces, affecting the rate of sea-level rise. \\

\noindent To proliferate in extreme environments, microbes have developed survival strategies, including dormancy. Dormancy is a reversible state of low metabolic activity utilised by microorganisms when energetic resources become scarce \cite{BradleyJamesA.2019Sotf}. Dormant microorganisms generate a seed bank of individuals, ready for resuscitation when favourable environmental conditions return \cite{JayT.Lennon2011Msbt}. \\

\noindent Dormancy is significant in microbial biogeography because it reduces dispersal limitation (stochastic mechanism) and resource competition (deterministic mechanism) \cite{locey2010synthesizing}. Dormancy must then be considered when quantifying the contributions of deterministic and stochastic mechanisms to microbial biogeography and its implication for a changing polar habitat. Despite this, little is known about the interplay between biogeography and dormancy.\\

\noindent Understanding the microbial ecology and biogeography of cryoenvironments is essential as both poles experience rapid climate change, affecting both soil and glacial systems \cite{SicilianoStevenD2014Sfia}. Historical records show that Arctic air temperature has warmed twice as fast as the global mean and may warm by more than 4\degree C by the end of the 21st Century \cite{stocker2014climate}. The adaptation of microorgamisms to extreme polar environments makes them among the most vulnerable to climate change \cite{WarwickFVincent2010Mert}. \noindent Previous research has shown that the recent rapid retreat of glaciers as a result of global warming exposes Arctic soils to microbial colonisation \cite{10.3389/feart.2017.00026}. As temperatures rise in the Arctic, microbial decomposition of the active layer and permafrost soil carbon could positively affect warming. \\ 

\noindent This project will quantify the activity and dormancy of microbes in Arctic glaciers and soils in a biogeographical framework. I will conduct fieldwork to investigate the spatial determinants of microbial community patterns. I will incorporate novel methods to access the active state of those communities. Finally I will implement this empirical data in developing a predictive numerical model, integrating biogeography and dormancy.\\

\chapter{Aims}

1. To quantify how geographic distance and habitat area affect microbial community diversity and dissimilarity in Arctic glaciers and what proportion of these are dormant.\\

\noindent 2. To quantify how geographic distance and habitat area affect microbial community diversity and dissimilarity in soils and what proportion of these are dormant.\\

\noindent 3. To quantify how these relationships change throughout the year.\\

\noindent 4. To develop a predictive numerical model synthesising microbial dormancy with biogeography in Arctic glaciers and soils.\\

\noindent 5. To use the numerical model to predict how microbial biogeography and dormancy change throughout the year and into the future as a result of climate change.

\chapter{Methods}

\section{Quantifying microbial dormancy and biogeography in Arctic glaciers and soils}

Previous investigations into microbial biogeography in the Arctic include bacteria in cryoconite holes \cite{darcy2018island}, red snow algae \cite{lutz2016biogeography} and large-scale investigations of soil diversity. \\

\noindent The general methodology employed in these investigations includes:\\
\noindent 1. Selecting samples sites, either with a spatially nested sampling scheme or far separated locations.\\
\noindent 2. Measuring habitat area and environmental properties (pH, conductivity, temperature, photosynthetic radiation, UV radiation, surface albedo, elevation, latitude, organic carbon) as well as GPS coordinates.\\
\noindent 3. Collecting samples and storing in Whirl-Pak bags or sterile test tubes before transporting to local facility for immediate DNA extraction using PowerSoil DNA extraction kit and/or storage.\\
\noindent 4. Transport to home research facility if applicable.\\
\noindent 5. PCR amplified 16s rRNA gene sequencing. \\
\noindent 6. Calculate diversity using Faith's phylogenetic diversity metric or another appropriate metric.\\
\noindent 7. Use spatial measurements (area/GPS coordinates) to assess any taxa-area relationships or compositional dissimilarity between communities.\\

\noindent *IMPORTANT NOTE* How can we develop ways to tackle the issue of unsampled rare taxa? \cite{woodcock2006taxa}

\section{Assessing active, dead and dormant community composition}

\noindent Whitfield 2005 Biogeography:  is  everything  everywhere? 117 additional species (for an original 20) found after incubating sample.\\

\noindent Metagenomics and 16S rRNA gene sequencing techniques fail to distinguish between living, dead and dormant cells. This is particularly significant for polar regions where preservation of dead cells and presence of viable dormant cells may skew our assessment of biodiversity and in situ activity. Bioorthogonal noncanonical amino acid tagging (BONCAT) is a high-throughput approach to detecting the incorporation of synthetic amino acids into proteins in active cells \cite{hatzenpichler2016visualizing}. BONCAT offers an alternative to methods that rely on expensive biomolecules or destructive techniques that prohibit downstream cell sorting and genomic sequencing (i.e. microautoradiography, secondary ion mass spectroscopy and Raman microspectroscopy). Other methods for quantifying microbial activity/dormancy include: \cite{burkert2019changes}:\\

\noindent1. LIVE/DEAD differential staining with microscopy\\
\noindent2. Endospore enrichment\\
\noindent3. Selective depletion of DNA\\

\noindent It might be useful to apply and assess the relative accuracy and practicality of a variety of activity/dormancy detection strategies. I will combine microbial biogeographic techniques (2.2) with cell state techniques (2.3) to explore both spatial and environmental patterns in microbial community structure, as well as the proportion of those communities in dead or dormant states. Data collection will take place over spring (April) and summer (August), allowing for a comparison of biogeographic patterns and state changes across initial and peak melt seasons. This data will help inform the predictive model in accounting for temperature change.

\subsection{Pole-to-Pole Connections: Similarities between Arctic and Antarctic Microbiomes and Their Vulnerability to Environmental Change \cite{kleinteich2017pole}} 

\noindent Using high-throughput 16S rRNA gene amplicon sequencing , terrestrial and lacustrine biofilms from the Arctic, Antarctic and temperate regions were characterised. Polar community compositions more similar to each other than to geographically closer temperate regions. Due to anthropogenic dispersal/environmental filtering?  \\

\section{Building a predictive numerical model}

\noindent I will use the empirical data collected to develop and test a predictive numerical model that will incorporate spatial and environmental determinants of microbial community state and structure. I will use this model to predict the changes to state and structure through seasonal variations and long term climate change. 


\subsection{Summary of Locey's Synthesizing traditional biogeography with microbial ecology: the importance of dormancy \cite{locey2010synthesizing}}

\noindent A conceptual model synthesising the ecological neutral theory of biogeography and microbial dormancy has previously been presented \cite{locey2010synthesizing} and may provide a base for model development.\\
{$\ast$ How important are dispersal limitations to microbes?}\\
{$\ast$ Increasing information about microbial spatial distribution.}\\
{$\ast$ Need more work on combining microbial ecology with their biogeography (i.e. how do their interactions with their environment, their unique microbial traits, affect their distributions and spatial patterns?)}\\
{$\ast$ Dormancy is a common, important microbial fitness trait.}\\

\noindent Locey considers dormancy as a biogeographic response. He looks at how dormancy might affect the equilibrium theory of island biogeography (how many species exist at equilibrium between colonisation and extinction rate) or on the unified neutral theory of biodiversity and biogeography (where each species within a habitat are considered ecologically equivalent and we do not need niche differences to explain patterns of biodiversity - these patterns arise from a combination of speciation, extinction and immigration from source populations \cite{harpole2010neutral}. \\

\noindent Locey also suggests that the equilibrium theory of island biogeography can produce patterns of microbial ubiquity, punctuated only by environmental selection (e.g. pH, salinity, temperature) in agreement with the Baas-Becking hypothesis (a niche assembly hypothesis) \cite{baas1934geobiologie}. Finally, Locey presents a conceptual model of the unified neutral theory biodiversity and biogeography with dormancy.\\

\noindent Dispersal limitation prevents non-porous range and defines the species geographic range. Dispersal limitation is important for evolution as a driver of speciation and extinction. Traditionally dispersal limitation unimportant to microbes due to resilient life stages, affected by environmental selection only \cite{baas1934geobiologie}. So we only need to understand the niche?  
{$\ast$ Difficulties in challenging this hypothesis come from differentiating dispersal limitation from environmental filtering. Small sample sizes or large communities also lead to rare taxa being underreported and biodiversity being under calculated \cite{woodcock2006taxa}. }\\
{$\ast$ Microbes thought to be able to disperse freely due to their small individual size and large population size \cite{fenchel2004ubiquity}.}\\
{$\ast$ <10\% of bacterial cells in terrestrial and aquatic environments may be active }\\

\noindent Forms of dormancy:\\
{$\ast$ Spores and cysts (physical)}\\
{$\ast$ Reduced metabolic rate/reproduction (physiological)}\\

\noindent Dormancy individuals can be easily reactivated and viable or cryptobiotic (devoid of all measurable metabolic activity) \\

\noindent Fundamental Biogeographic Responses:\\
{$\ast$ Adaptation}\\
{$\ast$ Dispersal (changing geographic range)}\\
{$\ast$ Extinction (local or global)}\\
{$\ast$ Dormancy? An alternative to extinction. Aids dispersal (although slow and uncertain), track environmental changes, maintain or expand range. }\\

\noindent Equilibrium Theory\\
{$\ast$ As more species colonise, less of the arriving individuals will represent new species, thus immigration rate (of new individuals) decreases as island species richness increases.}\\
{$\ast$ Higher species richness means more species can become extinct and so extinction rate rises}\\
{$\ast$ When immigration and extinction reach similar rates, species richnesses reaches equilibrium.}\\
{$\ast$ Immigration will be generally higher for closer islands and lower for isolated islands}\\
{$\ast$ Extinction will be generally greater for smaller islands and lower for larger islands}\\

\noindent Dormancy\\
{$\ast$ An alternative to extinction}\\
{$\ast$ Higher species richness = more extinction OR dormancy. This alternative will generally decrease extinction rates.}\\
{$\ast$ Dormancy = higher species richness at equilibrium (because new species arrive and persist in dormant states and do not compete to drive extinction)}\\
{$\ast$ If dormant individuals stay dormant they will eventually go extinct  unless maintained by immigrant individuals (as per the rescue effect)}\\
{$\ast$ As the source poor increases the number of species capable of dormancy, island communities will experience increasing species richness (like being on a large island close to the mainland).}\\
{$\ast$ Increasing species richness = more resource competition = more dormancy (less likely to reactivate)}\\
{$\ast$ Increasing extinction with increasing dormancy because populations cannot sustain themselves, they rely on immigrants (no reproduction)}\\
{$\ast$ Equilibrium theory only looks at total species richness, not relative abundance of species - we need ecological neutral theory! Need this to account for rare taxa. }\\
{$\ast$ Dormancy + ecological neutral theory incorporates stabilising mechanisms (storage effect) and equalising mechanisms (fitness equivalence).}\\
{$\ast$ Storage effect - why are there so many species in one community? As the environment changes, no species can be the best! They all respond differently, so they all store gains in different ways at different times.}\\
{$\ast$ UNT, can we explore microbial community dynamics in a neutral way (ecologically equivalent taxa) or do we need to take into account species idiosyncrasies. }\\

\noindent Incorporating Dormancy into Unified Neutral Theory\\
{$\ast$ Divide communities into active (no deaths?!) and dormant (no births) pools}\\
{$\ast$ Each individual becomes dormant before dying (Do microbes typically become dormant before dying?)}\\
{$\ast$ Disturbance drives individuals to dormancy, making room for reactivation into the active pool or reproduction by an already active individual. }\\
{$\ast$ Dormant pool non-competitive so it can grow in size}\\
{$\ast$ UNT active community is fixed by the total community can grow in size}\\
{$\ast$ What about incorporating horizontal gene transfer instead of just asexual reproduction?}\\

\subsection{Is everything everywhere? Stochastic and deterministic mechanisms in microbial biogeography (Solman, 2020)}
    
\noindent I have taken the model by Chisholm et al (2016) that describes a biphasic SAR (niche structured regime (low z value) to colonisation-extinction regime (higher z value)) and applied it to microbial SARs to see if a) there is a positive SAR b) is that SAR biphasic c) does the taxonomic group or habitat type affect the area at which that transition occurs? I.e. At low areas species richness should be determined by the number of niches and at a certain area immigration rate becomes more important in determining species richness. This happens at lower areas where immigration is high. So we think species that can immigrate more easily, and for less isolated habitats, the transition from niche-structures regime to colonisation extinction regime will be lower.\\

\noindent Many of my datasets did not have SARs suggesting that these are in a fulling deterministic, niche-structured regime (i.e. other factors may be important in determining species richness such as environmental filtering seen in archaea in soda lakes, or a lack of energy resource increase with area for fungi on leaves as expressed by the species-energy theory \cite{FeinsteinLarryM2012Tran}). Alternatively, there may be issues in sampling rare taxa \cite{woodcock2006taxa} or using inappropriate spatial scales \cite{davison2018microbial}.\\

% the z value is so small it mostly operates in a niche-structured regime or is it so ubiquitously high to be

\noindent Many datasets with positive SARs could not be successfully fit with the model, suggesting they opperate in an entirely colonisation-extinction driven regime. This was supported by the high immigration rates estimated by the model. The majority of these were aquatic and may be due to uncertainty in the spatial sample regime of a heterogenous habitat. In order to elucidate the the spatial patterns within these habitats, it might be useful to take a stratified approach. Some of the failed fittings had too few data points in comparison to the parameters of the model.\\

\noindent Of the successfully fitted datasets, habitat type did not have a significant affect on the area of transition between regimes (where a terrestrial island in marine waters or a biomembrane reactor might be considered more isolated than an inland lake). I believe this is due to too few data sets/too large a range in estimated critical areas for the datasets for any relationship to be found. Taxonomic group however, was significant. This may be because microbial taxa are more constrained by their physiology.ecology rather than habitat. They can overcome geographic barriers to dispersal. Bacteria have the lowest critical area suggesting they have higher dispersal rates/less dispersal limitation and species richness is a result of colonisation/extinction rates rather than habitat heterogenity. Bacteria can get there easily and aren't too bothered about the niches when they do. Pathogens showed the highest critical area, and I suggest this is due to their reliance on appropriate host species. The absence of host species may act as a major barrier to dispersal. It's not just a case of getting there, they also need the right niches to inhabit once they arrive.\\

\noindent So in conclusion, not all microbes exhibit positive SARs and this may be due to sampling error or environmental/niche factors being more important than colonisation extinction balance. \\

\noindent Not all positive SARs are biphasic. Some of them are entirely colonication-extinction based where environmental/niche factors do not play a significant role.\\

\noindent The data used in this analysis suggests taxonomic group and not habitat type are important in influencing at what point the transition between these regimes occurs. Taxa are constrained by their own physiology and biological needs and ability to overcome dispersal barriers, rather than the relative isolation of the habitats into which they disperse. 

\noindent I want to extend this work to look at the relative importance of deterministic and stochastic processes in the microbial biogeography of Arctic soils and glaciers, whilst incorporating Locey's conceptual model of dormancy.   \\           

\noindent There may also be an opportunity to incorporate the Microbial Ecophysiology in Low Energy Environments (MicroLow) model that considers the physiological transitions, varying energy requirements and flexible metabolic pathways of microorganisms through active and dormant states \cite{bradley2019survival}.  \\

\subsection{Widespread energy limitation to life in global subseafloor sediments \cite{bradley2020widespread}}

\noindent Novel mathematical model describing how organisms cope with extreme energy limitation over extended time periods, incorporating:\\
a) physiological transitions between active and dormant states\\
b) vary maintenance energy between these states\\
c) flexibility of energy source from exogenous (external origin) and endogenous (internal origin) catabolism (destructive metabolism)\\

\noindent {$\ast$ Microbial biomass divided into four pools (active, first stage of dormancy, second stage of dormancy, their stage of dormancy}\\
{$\ast$ Single pool of particulate organic carbon (POC) }\\
{$\ast$ Only the active pool can grow and growth rate is dependent on POC concentration  }\\
{$\ast$ Dormant cells are viable but needs to activate before it can grow}\\
{$\ast$ Over time cells transition to a lower state of dormancy }\\
{$\ast$ Cells from all pools can die, utilise POC for maintenance and transition to active or dormant stages}\\
{$\ast$ Deactivation/reactivation depends on supply of catabolic energy in relation to limiting factors such as POC}\\

\noindent The MicroLow model predicts biomass (number of individuals) and POC concentrations (amount of particulate organic carbon used up in an environment by cellular maintenance) in sediment over time. \\

\noindent \textbf{I want to incorporate the MicroLow model with the biogeographic conceptual model presented by Locey and the biphasic biogeographic model presented by Chisholm et al (2016) to try and predict the deterministic and stochastic mechanisms defining the spatial distribution of microbes in active and dormant states and their relative metabolic processes. I want to compare model simulations drawing from these three frameworks and empirical data collected through fieldwork. I hope to work with Anne and her team at the NHM to develop my skills in gene sequencing and techniques such as BONCAT for quantifying activity and dormancy.}


\chapter{Conclusion}

\noindent Dormancy in microbial systems helps to maintain biodiversity and stabilise ecosystem processes \cite{JayT.Lennon2011Msbt}. As Arctic permafrost becomes permanently degraded,  frozen soil organic matter is being made available to decomposers, thus releasing millions of tons of greenhouse gases \cite{JahnMarkus2010Gcca}. Through application of the predictive model, this project is relevant to decision making on climate change mitigation and conservation prioritisation. \\

\noindent Bioprospecting for culturable Arctic bacteria with cold-active enzymes has played an important part in developing industrial, agricultural and medical processes \cite{srinivas2009bacterial}. Arctic microbes have useful agricultural applications in regulating optimal fish health \cite{HamiltonErinF2019AACM} and in food processing such as milk fermentation and ice-cream production \cite{cid2016properties}. Through this investigation there may be scope for understanding new and bio-technologically useful species. \\

\noindent The Arctic contains some of the highest fidelity Mars analogue sites in the world \cite{pollard2009overview} and provides us with exciting opportunities to develop our understanding of interplanetary biology. A previous study has found indications of microbial metabolism in Arctic permafrost brine lenses formed during the Quaternary \cite{GilichinskyD2003Swbw}. These supercooled echoniches serve as a model for extraterrestrial life on cryogenic planets. By studying dormancy mechanisms we may glean insight into the life processes of extinct or extant extraterrestrial organisms. \\

\noindent Throughout this project I will engage with professional bodies to disseminate our findings. I will apply for the Royal Society of Biology's (RSB) Outreach and Engagement grant scheme to share our work with colleagues and the public. Moreover, I will present our findings during Biology Week 2021 and 2022. I will apply for the RSB Big Biology Day grant scheme, so that I can take our research into schools to engage children, particularly those from underrepresented groups. I will apply for the Microbiology Society's Conference Grant to share our work with colleagues. Moreover, I will submit articles detailing our research to Microbiology Today and other appropriate publications.\\

\noindent In conclusion, this project has far reaching implications for climate research, biotechnology and as a terrestrial analogue for extraterrestrial environments. I will disseminate our work and findings through formal and informal routes, engaging both colleagues and the general public. 

\bibliographystyle{apalike}
\bibliography{bibliography}
\addcontentsline{toc}{chapter}{Bibliography}
\end{document}